\documentclass[12pt]{article}
\usepackage{amsmath, amssymb}
\usepackage[margin=1in]{geometry}
\setlength{\parskip}{6pt}
\title{QUBO Formulation: Cyber Risk Scoring Node Selection Problem}
\date{}
\begin{document}
\maketitle

\subsection*{Cyber Risk Scoring Node Selection Problem}

\paragraph{Problem Statement.}
Given an infrastructure graph comprising $N$ assets, each asset $i$ is associated with a risk score $w_i$. Additionally, every distinct pair of assets $(i,j)$ is assigned a propagation-overlap penalty $p_{ij}$ that quantifies the redundancy or shared risk pathway between them. The objective is to select exactly $K$ assets for priority remediation or inspection such that the total risk-prioritization score is maximized. The total score is defined as the sum of the risk scores of the selected assets minus the sum of the pairwise overlap penalties among the selected assets.

\paragraph{Decision Variables.}
For each asset $i \in \{0, 1, \dots, N-1\}$, we introduce a binary decision variable $x_i \in \{0,1\}$. The variable $x_i$ equals $1$ if asset $i$ is selected for remediation or inspection, and $0$ otherwise.

\paragraph{QUBO Derivation.}
We convert the maximization problem into a minimization problem suitable for QUBO solvers by negating the objective function.

\textbf{Step 1 -- Original Objective.}
The original maximization objective is
\begin{align}
    \max_{x} \left( \sum_{i=0}^{N-1} w_i x_i - \sum_{0 \le i < j \le N-1} p_{ij} x_i x_j \right).
\end{align}
Negating yields the minimization objective:
\begin{align}
    \min_{x} \left( -\sum_{i=0}^{N-1} w_i x_i + \sum_{0 \le i < j \le N-1} p_{ij} x_i x_j \right).
\end{align}

\textbf{Step 2 -- Constraint.}
The requirement to select exactly $K$ assets is expressed as the equality constraint:
\begin{align}
    \sum_{i=0}^{N-1} x_i = K.
\end{align}

\textbf{Step 3 -- Penalty Term Expansion.}
We enforce the constraint by adding a quadratic penalty term $\lambda \left( \sum_{i=0}^{N-1} x_i - K \right)^2$ to the objective, where $\lambda > 0$ is a penalty weight. Expanding the square gives:
\begin{align}
    \lambda \left( \sum_{i=0}^{N-1} x_i - K \right)^2 
    &= \lambda \left( \sum_{i=0}^{N-1} x_i^2 + 2 \sum_{0 \le i < j \le N-1} x_i x_j - 2K \sum_{i=0}^{N-1} x_i + K^2 \right).
\end{align}
Using the idempotent property of binary variables, $x_i^2 = x_i$, we collect linear and quadratic terms:
\begin{align}
    \lambda \left( \sum_{i=0}^{N-1} x_i + 2 \sum_{0 \le i < j \le N-1} x_i x_j - 2K \sum_{i=0}^{N-1} x_i + K^2 \right)
    &= \sum_{i=0}^{N-1} \lambda(1 - 2K) x_i + \sum_{0 \le i < j \le N-1} 2\lambda x_i x_j + \lambda K^2.
\end{align}

\textbf{Step 4 -- Combined Hamiltonian.}
Adding the negated objective and the expanded penalty term yields the QUBO Hamiltonian $H(x)$:
\begin{align}
    H(x) = \sum_{i=0}^{N-1} \left( -w_i + \lambda(1 - 2K) \right) x_i + \sum_{0 \le i < j \le N-1} \left( p_{ij} + 2\lambda \right) x_i x_j + \lambda K^2.
\end{align}

\textbf{Step 5 -- Q Matrix and Offset.}
Reading off the coefficients from $H(x)$, the symmetric QUBO matrix entries $Q_{i,j}$ and constant offset $c$ are given by:
\begin{align}
    Q_{i,i} &= -w_i + \lambda(1 - 2K), & \text{for } i \in \{0, \dots, N-1\}, \\
    Q_{i,j} &= p_{ij} + 2\lambda, & \text{for } 0 \le i < j \le N-1, \\
    Q_{j,i} &= Q_{i,j}, & \text{for } 0 \le i < j \le N-1, \\
    c &= \lambda K^2.
\end{align}

\paragraph{Penalty Weight Justification.}
The penalty weight $\lambda$ must be chosen sufficiently large to ensure that any violation of the equality constraint is strictly penalized beyond the maximum possible improvement in the original objective. Since the maximum possible gain from the linear risk scores is bounded by $\max_{i} w_i$, setting $\lambda > \max_{i} w_i$ guarantees that the energy increase from violating $\sum_i x_i = K$ outweighs any potential reduction in the original objective terms. This ensures that the global minimum of $H(x)$ cannot reside in the infeasible region.

\paragraph{Correctness Argument.}
Minimizing $H(x)$ over $\{0,1\}^N$ recovers the optimal solution to the original problem because feasible assignments (those satisfying $\sum_i x_i = K$) incur zero penalty, making $H(x)$ identical to the negated original objective. For any infeasible assignment, the penalty term contributes at least $\lambda$ to the energy. Since $\lambda$ exceeds the maximum possible objective gain from any feasible selection, every infeasible assignment has strictly higher energy than the best feasible assignment. Consequently, the global minimum of $H(x)$ must correspond to a feasible assignment that minimizes the negated objective, which is equivalent to maximizing the original risk-prioritization score.

\paragraph{Illustrative Example.}
Consider a small instance with $N=3$ assets and $K=2$ required selections. Let the risk scores be $w_0 = 10$, $w_1 = 8$, and $w_2 = 12$. Let the pairwise overlap penalties be $p_{01} = 3$, $p_{02} = 5$, and $p_{12} = 2$. We choose $\lambda = 15$, which satisfies $\lambda > \max_i w_i$.

Substituting these values into the general formulas yields the concrete Q matrix entries:
\begin{align}
    Q_{0,0} &= -10 + 15(1 - 4) = -55, &
    Q_{1,1} &= -8 + 15(1 - 4) = -53, &
    Q_{2,2} &= -12 + 15(1 - 4) = -57, \\
    Q_{0,1} &= 3 + 2(15) = 33, &
    Q_{0,2} &= 5 + 2(15) = 35, &
    Q_{1,2} &= 2 + 2(15) = 32,
\end{align}
with offset $c = 15 \times 2^2 = 60$. The resulting Hamiltonian is:
\begin{align}
    H(x) = -55x_0 - 53x_1 - 57x_2 + 33x_0x_1 + 35x_0x_2 + 32x_1x_2 + 60.
\end{align}
Enumerating the three feasible assignments ($\sum_i x_i = 2$):
\begin{itemize}
    \item $x = (1,1,0)$: $H = -55 - 53 + 33 + 60 = -15$ (original score: $10+8-3=15$).
    \item $x = (1,0,1)$: $H = -55 - 57 + 35 + 60 = -17$ (original score: $10+12-5=17$).
    \item $x = (0,1,1)$: $H = -53 - 57 + 32 + 60 = -18$ (original score: $8+12-2=18$).
\end{itemize}
The global minimum is $H^* = -18$ at $x^* = (0,1,1)$, corresponding to the maximum risk-prioritization score of $18$. This confirms that assets $1$ and $2$ are optimally selected for remediation.

\end{document}